%=============================================================================!
% 15.7  RESULTS THAT DEPEND ON THE SIZE OF THE BOX                            !
%=============================================================================!
\section{Results that depend on the size of the box}
\label{sec:cv-size}

An error bar measures how far a run's average scatters about the average
that its own system would give in an endless run. It cannot measure a
bias that every run of that system shares, and a common such bias is the
size of the periodic box. This section measures the diffusion
coefficient of the liquid in boxes of 256 to 2048 atoms, finds it changes
by more than any error bar, and accounts for the change.

\subsection*{A ball in a narrow tube}

A ball sinking through honey falls more slowly in a narrow tube than in a
wide pot. The honey it pushes down has to flow back up past it, through
the gap between the ball and the wall, and the narrower the gap the
faster that return flow and the stronger its drag. A periodic box has no
walls, but it has the same return flow. An atom moving through the liquid
drags its surroundings with it, and in a periodic box every image of the
atom drags the same flow in every image of the box; since the liquid as a
whole does not move, the fluid must flow back between them, against the
atom, and the more so the closer the images are.

How strongly a fluid resists such flows is its shear viscosity
$\eta_{\mathrm s}$: a layer of fluid of thickness $d$ between two plates,
one moving past the other at the speed $u$, needs the force $\eta_{\mathrm
s}u/d$ per unit area of plate to keep it moving, so $\eta_{\mathrm s}$ is in
\si{\pascal\second}. Honey has a large $\eta_{\mathrm s}$, water a small
one. The only combination of $\kB T$, $\eta_{\mathrm s}$ and the side $L$
of the box that has the units of a diffusion coefficient is $\kB
T/(\eta_{\mathrm s}L)$ (\cref{ex:cv-units}), so the return flow should
lower $D$ by a number times $\kB T/(\eta_{\mathrm s}L)$. Summing the flows
of all the images of a cubic box gives the number, and Yeh and Hummer
showed that molecular dynamics follows
\begin{equation}
   D(L) = D_\infty - \frac{2.837297\,\kB T}{6\pi\eta_{\mathrm s}L} ,
   \label{eq:cv-yh}
\end{equation}
with $D_\infty$ the coefficient of an endless liquid\cite{Yeh2004,
Dunweg1993}. The correction falls only as $1/L$, that is as $N^{-1/3}$ at
a fixed density: eight times the atoms halve it.

\subsection*{Five boxes}

The liquid at its density of $0.8\sigma^{-3}$ and \SI{135}{\kelvin} was
prepared as in Chapter~12 with 256, 500, 864, 1372 and 2048 atoms, boxes of
side 23.26 to \SI{46.51}{\angstrom}, and followed under CSVR for
\SI{1}{\nano\second} from each of two seeds. $D$ comes from twenty blocks
of \SI{100}{\pico\second} at each size, each moved to \SI{135}{\kelvin}
along the slope $\dd D/\dd T$ of \cref{sec:cv-observables}, a correction
of at most \SI{0.009e-4}{\angstrom\squared\per\femto\second}. In units of
\SI{e-4}{\angstrom\squared\per\femto\second} it rises with the box: $3.914
\pm 0.032$, $4.062 \pm 0.017$, $4.166 \pm 0.022$, $4.200 \pm 0.015$ and
$4.214 \pm 0.011$ (\cref{fig:cv-size}(a)). The change from the smallest box
to the largest is nine times the larger of their two errors. A straight
line in $1/L$, fitted as in \cref{sec:cv-drift} but with each squared miss
divided by the square of that box's error, so that the precise boxes count
for more, has $D_\infty = 4.503 \pm 0.032$ in the same units and the slope
$-12.88 \pm 1.19$ in \SI{e-4}{\angstrom\cubed\per\femto\second}. Their
errors follow as in the toolbox of \cref{sec:cv-drift}: with the weights
$1/\sigma_i^2$, $\sigma_i$ the error of box $i$, the slope is a weighted sum
of the $D$ of the boxes whose variance, each term's weight squared times
$\sigma_i^2$, adds up to one over $\sum_i(x_i - \bar x)^2/\sigma_i^2$, with
$x_i = 1/L_i$ and $\bar x$ their weighted mean. The sum
of the squared misses in units of the errors is 5.3; each such term
averages 1 with the spread $\sqrt2$ (\cref{ex:cv-error-error}), and
fitting two numbers to five boxes leaves three, so scatter alone gives $3
\pm 2.4$. The box of 256 atoms used through
Chapters~12 to 14 gives a $D$ 13.1\% below that of the endless liquid, and
even 2048 atoms 6.4\% below. Leaving out the smallest box moves
$D_\infty$ to 4.470, by about its error.

Through \cref{eq:cv-yh}, the slope gives the viscosity
\begin{equation*}
   \eta_{\mathrm s} = \frac{2.837297\,\kB T}{6\pi\,\lvert\text{slope}\rvert}
   = (2.18 \pm 0.20)\times10^{-4}\,\si{\pascal\second} ,
\end{equation*}
or $2.40 \pm 0.22$ in the reduced unit $\sqrt{m\varepsilon}/\sigma^2 =
\SI{9.07e-5}{\pascal\second}$ of the Lennard-Jones model. Chapter~16 measures $\eta_{\mathrm s}$ independently,
from the fluctuations of the shear stress, which tests whether the whole
of the change with size is the return flow of \cref{eq:cv-yh}.

The static averages change far less. The potential energy per atom and
the pressure of each box lie within 2.1 combined errors of those for 2048
atoms (\cref{fig:cv-size}(b)). The largest difference, the
pressure of 256 atoms, is \SI{0.0009}{\giga\pascal}, 1\% of the
pressure, where $D$ differs by 7\% between the same two boxes: the small
box serves for this liquid's energy and pressure, and not for its
diffusion.

\begin{figure}[htb]
\centering
\includegraphics{ch15_convergence/size}
\caption{A result that depends on the size of the box: liquid argon at
\SI{135}{\kelvin} and $0.8\sigma^{-3}$ with 256 to 2048 atoms, two runs of
\SI{1}{\nano\second} at each size. (a) The diffusion coefficient, with its
error from blocks of \SI{100}{\pico\second}, against $1/L$, with the line
of \cref{eq:cv-yh} fitted to them (dashed) and its value at $1/L = 0$
(ochre). (b) The potential energy per atom and the pressure for each box
less those for 2048 atoms, in units of the error of the difference,
against $\pm2$ (dotted).}
\label{fig:cv-size}
\end{figure}

\begin{keyidea}
A converged average can still depend on the size of the box. The
diffusion coefficient of a periodic liquid falls short of the endless
liquid's by $2.837\kB T/(6\pi\eta_{\mathrm s}L)$, as the image flows return
through the box; runs at several sizes and a straight line in $1/L$ find
$D_\infty$. Static averages of the same liquid change far less.
\end{keyidea}

\notebook{15}{fit the line through the five boxes and extrapolate}

\begin{selfcheck}
By how much should $D$ change from a box of 500 atoms to one of 4000 at
the same density, if it follows the line of this section?
\end{selfcheck}
